% !TeX root = ../main.tex
% ---------------------------------------------------------------------------
% Rock Reinforcement and Bolting
% Scope extension: expands the support material in Chapters 6 and 10.
% Key sources: Windsor & Thompson on reinforcement classification; Hoek et al.
% Support of Underground Excavations in Hard Rock; Li on energy-absorbing
% bolts; AS 4671 and ISRM Suggested Methods for bolt testing.
% ---------------------------------------------------------------------------
\chapter{Rock Reinforcement and Bolting}
\label{ch:rock-reinforcement}


\section{Reinforcement principles}
\label{sec:reinforcement-principles}

\subsection{Reinforcement versus support}
\label{sec:reinforcement-vs-support}
\scopenote{Reinforcement acts within the rock mass to improve its own
capacity; support acts on the excavation boundary. The distinction determines
how each is analysed.}

\subsection{Classification of reinforcement systems}
\label{sec:reinforcement-classification}
\scopenote{The Windsor \& Thompson framework: continuously mechanically
coupled, continuously frictionally coupled and discretely mechanically
coupled devices, and what each implies for load transfer.}

\section{Bolt types and load transfer}
\label{sec:bolt-types-and-load-transfer}

\subsection{Bolt types and their mechanics}
\label{sec:bolt-types}
\scopenote{Mechanical anchors, fully grouted rebar, resin-anchored bolts,
friction bolts (split set, Swellex), self-drilling anchors and cable bolts.}

\subsection{Load transfer and bond behaviour}
\label{sec:load-transfer}
\scopenote{Shear stress distribution along the bond, decoupling and
progressive debonding, critical embedment length, and the effect of borehole
and grout quality.}

\section{Design for the governing mechanism}
\label{sec:design-for-the-governing-mechanism}

\subsection{Bolting for structurally controlled failure}
\label{sec:bolting-wedges}
\scopenote{Wedge and block stabilisation, required bolt force from limit
equilibrium, orientation relative to the sliding direction identified in
Chapter~\ref{ch:structural-analysis}, and factor of safety conventions.}

\subsection{Bolting for stress-induced failure}
\label{sec:bolting-stress}
\scopenote{Confinement provided by systematic bolting, its effect on the
ground reaction curve of Chapter~\ref{ch:ccm}, and retention of fractured rock
in spalling ground.}

\subsection{Pattern bolting design}
\label{sec:pattern-bolting}
\scopenote{Length and spacing rules from span and $Q$; empirical charts;
consistency checks between empirical, analytical and numerical bolt designs.}

\section{Shear, cables and dynamic reinforcement}
\label{sec:shear-cables-and-dynamic-reinforcement}

\subsection{Shear behaviour and dowel action}
\label{sec:dowel-action}
\scopenote{Bolts crossing a shearing discontinuity, combined axial and shear
loading, and the contribution of dowel action to joint shear resistance.}

\subsection{Cable bolts and long tendons}
\label{sec:cables}
\scopenote{Plain and bulbed strand, pre-tensioning, multi-strand anchors, and
their role in large-span caverns (Chapter~\ref{ch:caverns}).}

\subsection{Dynamic and yielding reinforcement}
\label{sec:dynamic-bolts}
\scopenote{Energy-absorbing bolts for burst-prone ground, energy demand and
capacity matching, and system compatibility with mesh and shotcrete.
Chapter~\ref{ch:squeezing-rockburst}.}

\section{Surface retention and durability}
\label{sec:surface-retention-and-durability}

\subsection{Surface retention: mesh, straps and plates}
\label{sec:surface-retention}
\scopenote{The often-neglected third element of a support system; mesh types,
plate stiffness and the importance of connection detailing.}

\subsection{Corrosion protection and design life}
\label{sec:bolt-corrosion}
\scopenote{Aggressivity classification, galvanising, encapsulation, double
corrosion protection, and the distinction between temporary and permanent
reinforcement.}

\section{Testing and quality control}
\label{sec:bolt-testing}
\scopenote{Pull-out and proof load testing, grout quality verification,
installation records, and typical acceptance criteria.}
